\section*{Chapter \ref{ch:b2:magmatic-series} --- Magmatic Series and Tectonic Settings}

\begin{solution}{exo:b2:magmatic-series:1}
(a) $\ce{SiO2} = 49.60$, alkalis $2.22 + 0.51 = 2.73$: between 45 and 52 and below 5,
\emph{basalt} (B). (b) 58.5 and $3.9 + 1.6 = 5.5$: between 57 and 63; the upper
boundary there runs from (57, 5.9) to (63, 7.0), at $5.9 + 1.1\times1.5/6 = 6.18$
for 58.5; 5.5 lies below it: \emph{andesite} (O2). (c) 72.0 and $4.0 + 4.2 = 8.2$:
beyond the vertical at 69 above 8 and above the line from (69, 8) to (77, 0), which is
at 5.0 for 72: \emph{rhyolite} (R). (d) 47.0 and $4.5 + 2.5 = 7.0$: at 47 the field U1
lies between the line from (45, 5) to (49.4, 7.3), at $5 + 2.3\times2/4.4 = 6.05$, and
the line from (45, 9.4) to (49.4, 7.3), at $9.4 - 2.1\times2/4.4 = 8.45$. The rock is
in U1: a \emph{tephrite} if its norm holds less than \qty{10}{\percent} olivine, a
\emph{basanite} if more.
\end{solution}

\begin{solution}{exo:b2:magmatic-series:2}
(a) The oxides add up to 98.40 and the total with volatiles to 100.00;
\ce{H2O} $1.40 < 2$ and \ce{CO2} $0.20 < 0.5$: the analysis is acceptable. (b) Factor
$100/98.40 = 1.01626$: \ce{SiO2} 52.44, \ce{TiO2} 1.22, \ce{Al2O3} 16.36,
\ce{FeO_T} 9.96, \ce{MgO} 6.71, \ce{CaO} 8.33, \ce{Na2O} 3.56, \ce{K2O} 1.42
(sum 100.00); alkalis $3.56 + 1.42 = 4.98$. (c) Before: 51.60 and 4.90, just below the
basalt--basaltic andesite boundary at 52: \emph{basalt}. After: 52.44 and 4.98; at
52.44 the upper boundary of O1, from (52, 5) to (57, 5.9), is at $5 + 0.9\times0.44/5
= 5.08$: \emph{basaltic andesite} (O1). The recalculation moved the rock across a
boundary: names near a line depend on it. (d) At 52.44 the low-K line is at $0.3 +
0.045\times4.44 = 0.50$ and the high-K line at $1.2 + 0.095\times4.44 = 1.62$; 1.42
lies between: \emph{medium-K}.
\end{solution}

\begin{solution}{exo:b2:magmatic-series:3}
(a) $Q + A + P = 90$: $Q = 25/90 = \qty{27.8}{\percent}$ and $100\,P/(A + P) =
100\times35/65 = 53.8$: field 3b, a \emph{monzogranite}, a granite. Biotite and
hornblende do not enter the QAPF sum. (b) $Q + A + P = 65$: $Q = 5/65 =
\qty{7.7}{\percent}$ (between 5 and 20) and $100\times55/60 = 91.7$ (above 90): field
10*, quartz diorite, quartz gabbro or quartz anorthosite. The \omterm{def:b2:magmatic-series:qapf}{QAPF diagram} alone does
not choose among them; the plagioclase and the mafic minerals (here
\qty{35}{\percent} clinopyroxene) do.
\end{solution}

\begin{solution}{exo:b2:magmatic-series:4}
The low-K line is $0.3 + 0.045(\ce{SiO2} - 48)$, the high-K line $1.2 +
0.095(\ce{SiO2} - 48)$. (a) At 50: 0.39 and 1.39; $0.40$ is just above the low line:
\emph{medium-K}. (b) At 56: 0.66 and 1.96: \emph{medium-K}. (c) At 62: 0.93 and 2.53;
$2.90 > 2.53$: \emph{high-K}. (d) At 66: 1.11 and 2.91; $1.00 < 1.11$: \emph{low-K}.
\end{solution}

\begin{solution}{exo:b2:magmatic-series:5}
(a) Basalt: $A + F + M = 3.2 + 14.5 + 4.5 = 22.2$; $A = 14.4$, $F = 65.3$, $M = 20.3$.
Andesite: $5.05 + 6.1 + 3.8 = 14.95$; $A = 33.8$, $F = 40.8$, $M = 25.4$. (b) The line
never reaches beyond $F = 54.5$; the basalt, with $F = 65.3$, lies on the side of F
whatever its $A$: \emph{tholeiitic}, a ferrobasalt of a Fenner trend. For the
andesite, at $A = 33.8$ the line lies between (49.6, 35.0) and (52.5, 30.0), at $F =
49.6 + 2.9\times1.2/5 = 50.3$; $40.8 < 50.3$: \emph{calc-alkaline}. (c) Magnetite,
which carries ferric iron out of an oxidised \omterm{def:b2:magma-emplacement:magma}{magma} from early on and keeps $F/M$ from
rising as fast as in a tholeiite.
\end{solution}

\begin{solution}{exo:b2:magmatic-series:6}
MORB: $\mathrm{Nb_N} = 4/0.658 = 6.08$, $\mathrm{La_N} = 4.12/0.648 = 6.36$, ratio
0.956. Ocean-island basalt: $38/0.658 = 57.8$, $30.8/0.648 = 47.5$, ratio 1.215. Arc
basalt: $3/0.658 = 4.56$, $6.1/0.648 = 9.41$, ratio 0.484. Only the arc basalt has a
niobium anomaly, a trough at Nb on a spider diagram: niobium is depleted relative to
its neighbour La: the arc \omterm{def:b2:magma-emplacement:magma}{magma} received elements from the slab, but little
niobium. MORB is nearly flat; the ocean-island basalt is richer
in Nb than in La.
\end{solution}

\begin{solution}{exo:b2:magmatic-series:7}
(a) $\mathrm{Ti} = 1.6\times0.5993\times10^4 = \qty{9589}{ppm}$; Ti/100 = 95.9, Zr
110, $3\mathrm Y = 96$; sum 301.9; coordinates 31.8, 36.4, 31.8: field B. (b) Ti
16\,780, so 167.8, 220, 87; sum 474.8; 35.3, 46.3, 18.3: field D. (c) Ti 4794, so
47.9, 50, 60; sum 157.9; 30.3, 31.7, 38.0: field B. Basalt (b) is a within-plate basalt
(of an ocean island or a continent); (a) and (c), in field B, may be MORB, arc
tholeiites or calc-alkali basalts: the diagram cannot decide, and another one (a
spider diagram, the Nb anomaly) is needed.
\end{solution}

\begin{solution}{exo:b2:magmatic-series:8}
(a) MORB $142/164 = 0.866$, ocean-island basalts $4/173 = 0.023$, arc basalts $101/201 =
0.502$. (b) With equal priors, Bayes' rule gives $0.866/(0.866 + 0.023 + 0.502) =
0.866/1.391 = 0.623$. (c) That the basalt is very unlikely to be an ocean-island basalt
(probability $0.023/1.391 = 0.017$), but not whether it is a MORB or an arc basalt:
roughly three chances in five against two.
\end{solution}

\begin{solution}{exo:b2:magmatic-series:9}
Moles per \qty{100}{g}: \ce{Na2O} $5.5/61.98 = 0.08874$, \ce{K2O} $4.5/94.20 =
0.04777$, \ce{Al2O3} $10.5/101.96 = 0.10298$. (a) $(0.08874 + 0.04777)/0.10298 =
0.13651/0.10298 = 1.326 > 1$: the rhyolite is \emph{peralkaline}. (b) Feldspar can
take $0.10298$ of the $0.13651$ moles of alkali oxide; the excess, $0.03353$, is
$0.03353/0.13651 = 0.246$, a quarter of the alkalis, which must enter other minerals.
\end{solution}

\begin{solution}{exo:b2:magmatic-series:10}
(a) $s = (\lambda_F - \lambda_M)/(\lambda_F - \lambda_A)$: for $\lambda_F = 0.5$, $(0.5
- 0.8)/(0.5 - 0.4) = -3.0$; for $0.7$, $(-0.1)/(0.3) = -0.333$. (b) Solving
$s(\lambda_F - 0.4) = \lambda_F - 0.8$ gives $\lambda_F = (0.4s - 0.8)/(s - 1)$:
Iceland, $s = -1.10$: $(-0.44 - 0.8)/(-2.10) = 0.590$; Cascades, $s = -0.72$:
$(-0.288 - 0.8)/(-1.72) = 0.633$. (c) The Cascades \omterm{def:b2:magma-emplacement:magma}{magma}, with the larger
$\lambda_F$, loses its iron faster: magnetite takes it out. (d) $\ln(F/M)$ grows at the
rate $\lambda_M - \lambda_F$ and $\ln(A/F)$ at the rate $\lambda_F - \lambda_A$; so for
each unit gained by $\ln(A/F)$, $\ln(F/M)$ gains $-s$: 1.10 in Iceland, 0.72 in the
Cascades. $F/M$ rises along both trends, but faster in the tholeiitic one, the Fenner
trend.
\end{solution}

\begin{solution}{exo:b2:magmatic-series:11}
(a) BHVO-2, moles per \qty{100}{g}: \ce{SiO2} 0.82557, \ce{Al2O3} 0.13182, \ce{FeO}
0.15521, \ce{MgO} 0.18007, \ce{CaO} 0.20328, \ce{Na2O} 0.03580, \ce{K2O} 0.00545.
Orthoclase 0.00545 and albite 0.03580; \ce{Al2O3} left $0.13182 - 0.04125 = 0.09057$,
all into anorthite with as much \ce{CaO}; \ce{CaO} left $0.11271$ into diopside;
$(\ce{Mg,Fe})\ce{O}$ left $0.33528 - 0.11271 = 0.22257$ into hypersthene. Silica
needed: $6\times0.00545 + 6\times0.03580 + 2\times0.09057 + 2\times0.11271 + 0.22257 =
0.87663$; available 0.82557; short by 0.05106. Each $(\ce{Mg,Fe})\ce{O}$ turned from
hypersthene into olivine saves half a silica, at most $0.22257/2 = 0.11129$: turning
$0.10212$ of it is enough, and $0.12045$ stays in hypersthene. The norm holds olivine
and hypersthene, and neither quartz nor nepheline: the rock is silica-saturated in the
sense of the definition, a hypersthene-normative \emph{olivine tholeiite}, of the
\omterm{def:b2:magmatic-series:series}{tholeiitic series}, as the shield stage of a Hawaiian volcano erupts. (b) Moles: 0.75732, 0.14712, 0.16008, 0.21092, 0.17832, 0.06131,
0.01592. Or 0.01592, ab 0.06131, an 0.06988, di 0.10843, hy $0.37100 - 0.10843 =
0.26257$; silica needed $0.09552 + 0.36786 + 0.13976 + 0.21686 + 0.26257 = 1.08257$;
available 0.75732; short by 0.32525. All the hypersthene turned into olivine saves
0.13129, leaving 0.19396; each mole of \ce{Na2O} turned from albite into nepheline
frees 4 \ce{SiO2}, so $0.19396/4 = 0.04849$
of the 0.06131 moles of albite \ce{Na2O} become nepheline. Nepheline in the norm: the
rock is silica-undersaturated, an \emph{alkali basalt}, of the \omterm{def:b2:magmatic-series:series}{alkaline series}.
\end{solution}

\begin{solution}{exo:b2:magmatic-series:12}
(a) Arc basalts $116/141 = 0.823$, MORB $21/101 = 0.208$, ocean-island basalts
$3/137 = 0.022$. (b) $0.823/(0.823 + 0.208 + 0.022) = 0.823/1.053 = 0.782$.
(c) Numerator $0.1\times0.823 = 0.0823$; denominator $0.0823 + 0.6\times0.208 +
0.3\times0.022 = 0.0823 + 0.1248 + 0.0066 = 0.2137$; probability $0.385$. (d) The
same test on the same rock gives 0.78 or 0.39 depending on how common each setting
is where the rock was found: a \omterm{def:b2:magmatic-series:discrimination}{discrimination diagram} gives a likelihood, and the
geological context must supply the rest.
\end{solution}

\begin{solution}{pb:b2:magmatic-series:1}
\noindent\textbf{Part I.}

\textbf{1.} Oxides: lava 1 99.17, lava 2 98.65, lava 3 98.20; with the volatiles
99.97, 99.80 and 99.45, all close to 100. Water (0.70, 1.10, 1.20) is below
\qty{2}{\percent} and carbon dioxide (0.10, 0.05, 0.05) below \qty{0.5}{\percent}: all
three are acceptable.

\textbf{2.} Factor $100/99.17 = 1.00837$: \ce{SiO2} 50.82, \ce{Na2O} 2.92, \ce{K2O}
0.63, alkalis 3.55.

\textbf{3.} Lava 2: $100/98.65 = 1.01368$, \ce{SiO2} 59.40, \ce{Na2O} 3.95, \ce{K2O}
1.17, alkalis 5.12. Lava 3: $100/98.20 = 1.01833$, \ce{SiO2} 66.50, \ce{Na2O} 4.48,
\ce{K2O} 1.88, alkalis 6.36.

\textbf{4.} Lava 1 (50.82, 3.55): basalt. Lava 2 (59.40, 5.12; the O2 upper boundary
there is at $5.9 + 1.1\times2.40/6 = 6.34$): andesite. Lava 3 (66.50, 6.36; between 63
and the line from (69, 8) to (77, 0), below the line from (63, 7) to (69, 8)): dacite.

\textbf{5.} Raw: (50.4, 3.52), (58.6, 5.05), (65.3, 6.25): the same fields. The
recalculation moves each point by at most \qty{1.2}{\percent} silica and changes
nothing here, unlike \cref{exo:b2:magmatic-series:2}.

\textbf{6.} The \omterm{def:b2:magmatic-series:qapf}{QAPF diagram}, from the measured mode. A plutonic rock of andesitic
composition is rich in plagioclase, with little alkali feldspar or quartz: it should
fall in field 10 or 10*, a diorite or quartz diorite (its mode decides).

\noindent\textbf{Part II.}

\textbf{7.} Lava 1: lines at 0.43 and 1.47, $\ce{K2O} = 0.63$: medium-K. Lava 2: 0.81
and 2.28, 1.17: medium-K. Lava 3: 1.13 and 2.96, 1.88: medium-K.

\textbf{8.} Volatile-free \ce{FeO_T} and \ce{MgO}: 8.67 and 7.97; 6.18 and 3.85; 4.07
and 1.63. Lava 1: $3.55 + 8.67 + 7.97 = 20.19$; $A = 17.6$, $F = 42.9$, $M = 39.5$.
Lava 2: 15.15; 33.8, 40.8, 25.4. Lava 3: 12.06; 52.7, 33.7, 13.5.

\textbf{9.} Lava 1, $A = 17.6$: the line between (42.8, 17.1) and (46.9, 18.1) is at
$F = 42.8 + 4.1\times0.5/1.0 = 44.9$; $42.9 < 44.9$, calc-alkaline side, close to the
line. Lava 2, $A = 33.8$: $F = 49.6 + 2.9\times1.2/5 = 50.3$; $40.8 < 50.3$,
calc-alkaline. Lava 3, $A = 52.7$: $F = 40.8 - 2.9\times2.7/5 = 39.2$; $33.7 < 39.2$,
calc-alkaline.

\textbf{10.} $\ce{FeO_T}/\ce{MgO}$: $8.67/7.97 = 1.09$, $6.18/3.85 = 1.61$, $4.07/1.63 =
2.50$. The ratio rises, as it does in every differentiating basalt, but iron itself
falls from 8.67 to 4.07, and $F$ falls on the triangle from 42.9 to 33.7: no iron
enrichment. A tholeiitic suite would first climb towards F, as the ferrobasalt of
\cref{exo:b2:magmatic-series:5} does.

\textbf{11.} The \omterm{def:b2:magmatic-series:series}{calc-alkaline series}. Magnetite, crystallising from an oxidised \omterm{def:b2:magma-emplacement:magma}{magma},
removed iron with magnesium; the oxidised species come from the dehydrating slab of a
subduction zone.

\noindent\textbf{Part III.}

\textbf{12.} $\mathrm{Nb_N} = 2.0/0.658 = 3.04$, $\mathrm{La_N} = 7/0.648 = 10.80$,
$\mathrm{Th_N} = 1.1/0.0795 = 13.84$.

\textbf{13.} $3.04/10.80 = 0.281$, well below the medians of MORB (0.956) and
ocean-island basalts (1.215), and below even the median arc basalt (0.484): a deep
niobium trough.

\textbf{14.} $13.84/3.04 = 4.55$; the median MORB of the spider diagram has $2.64/6.08 =
0.43$. Thorium, like the other elements enriched in arc basalts, was added by the
subducting slab.

\textbf{15.} \ce{TiO2} volatile-free $0.95\times1.00837 = 0.958$; $\mathrm{Ti} =
0.958\times5993 = \qty{5741}{ppm}$; Ti/100 = 57.4, Zr 70, $3\mathrm Y = 54$, sum 181.4:
31.6, 38.6, 29.8, in field B, shared by MORB, arc tholeiites and calc-alkali basalts.
It excludes a within-plate basalt but does not decide between a ridge and an arc. The
diagram was built for basalts: in an andesite and a dacite, titanium has been removed
by the oxides and zirconium concentrated by crystallisation, and their fields (lava 2
would fall in C, lava 3 outside every field) would mean nothing.

\textbf{16.} A calc-alkaline, medium-K series, with a deep niobium trough and high
$\mathrm{Th_N}/\mathrm{Nb_N}$, on a volcano that erupted basalt to dacite: the volcano
grew in a volcanic arc above a subduction zone.

\noindent\textbf{Part IV.}

\textbf{17.} $C_L = C_0/F$ for a perfectly \omterm{def:b2:elements-in-the-earth:partition}{incompatible element}: $F = 0.63/1.17 =
0.538$.

\textbf{18.} Rubidium $10/20 = 0.500$; thorium $1.1/2.1 = 0.524$. (The \omterm{def:b2:elements-in-the-earth:trace}{trace
elements}, given as analysed, would change by less than \qty{1}{\percent} with the
recalculation.)

\textbf{19.} $C_L/C_0 = F^{D - 1}$: $125/70 = 1.786 = F^{-0.9}$, $F = 1.786^{-1/0.9} =
0.525$.

\textbf{20.} From potash, $F = 1.17/1.88 = 0.622$ from lava 2 to lava 3. Strontium:
$380/480 = 0.792 = 0.622^{D - 1}$, so $D - 1 = \ln 0.792/\ln 0.622 = -0.234/-0.474 =
0.49$ and $D \approx 1.49$: strontium has become compatible. Plagioclase, which takes
strontium in, was crystallising in abundance.

\textbf{21.} From potash, $F = 0.63/1.88 = 0.335$ (rubidium $10/32 = 0.31$, thorium
$1.1/3.4 = 0.32$): about two thirds of the basalt, $1 - 0.335 = 0.665$, crystallised
before the dacite.

\textbf{22.} Assimilated potash would raise \ce{K2O} in lava 2 above what
crystallisation alone gives, so $F = 0.63/1.17$ would come out too low, and the mass of
crystals too high.

\textbf{23.} $(0.538 + 0.500 + 0.524 + 0.525)/4 = 2.087/4 = 0.522$; crystals $1 - 0.522
= 0.478$ of the basalt, that is \textbf{\qty[text-series-to-math=true]{48}{kg} of
crystals per \qty{100}{kg} of basalt} to make the andesite.
\end{solution}
